\chapter{LAUNCH CAMPAIGN PREPARATIONS}
\label{sec:campaign_preparations}
\section{Input for the Campaign / Flight Requirement Plan}
\label{sec:input_FRP}

%%%%%%%%%%%%%%%%%%%%%%%%%%%%%%%%%%%%%%%%%%%%%%%%%%%%%%%%%%%%%%%%%%%%%%%%%%%%%%%
\subsection{Dimensions and Mass}

\begin{table}[htbp]
\centering
\caption{Experiment Mass and Volume}
\label{tab:mass_volume}
\begin{tabular}{L{6cm}L{6cm}}
\toprule
Experiment mass (in kg): & $\sim$4\\
Experiment footprint area (for BEXUS) (in $m^2$): & 0.04\\
Experiment volume (in $m^3$): & 0.006\\
Experiment expected \gls{cog} position: & Distance from geometric centre: x=4mm y=23mm z=22mm \\
\bottomrule
\end{tabular}
\end{table}
\clearpage 

%%%%%%%%%%%%%%%%%%%%%%%%%%%%%%%%%%%%%%%%%%%%%%%%%%%%%%%%%%%%%%%%%%%%%%%%%%%%%%%
\subsection{Ground and Flight Safety Risks}

\begin{table}[htbp]
\centering
\caption{Ground and flight safety risks}
\label{tab:safety_risks}
\begin{tabular}{p{4cm} p{5cm} p{4cm}}
\toprule
\textbf{Risk} & \textbf{Key Characteristics} & \textbf{Mitigation} \\
\midrule
Use of a pressurised vessel with 3 bar difference beween outside and inside. & Material: Aluminium & Pressure test will be conducted up to 4 bar in line with the Esrange safety manual. \\
\bottomrule
\end{tabular}
\end{table}
\clearpage

%%%%%%%%%%%%%%%%%%%%%%%%%%%%%%%%%%%%%%%%%%%%%%%%%%%%%%%%%%%%%%%%%%%%%%%%%%%%%%%
\subsection{Electrical Interfaces}

\begin{table}[htbp]
\centering
\caption{Electrical interfaces applicable to BEXUS}
\label{tab:electrical_interfaces_bexus}
\begin{tabular}{p{10cm} p{5cm}}
\toprule
\multicolumn{2}{l}{\textbf{BEXUS Electrical Interfaces}} \\
\midrule
E-Link Interface: E-Link required? Yes/No & Yes\\
Number of E-Link interfaces: & 1 \\
Number of required IP addresses: & 2 (one for the experiment and one for the ground station)\\
Data rate -- downlink (max. and average): & max: 100 Kbit/s  (average: 33 Kbit/s) \\
Data rate -- uplink (max. and average): & max: 10 Kbit/s \\
Interface type (RS-232, Ethernet): & Ethernet \\
\midrule
Power system: Gondola power required? Yes/No & Yes\\
Peak power and current consumption: & 65.05 W and 2.26 A \\
Average power and current consumption: & 23.11 W and 0.8 A \\
Total power and current consumption after lift-off & 138.67 Wh and 8.59 Ah \\
\midrule
Power system: Experiment includes batteries? Yes/No & No\\
Type of batteries: & \\
Number of batteries (and S x P): & \\
Capacity (1 battery): & Ah \\
Voltage (1 battery): & V \\
\bottomrule
\end{tabular}
\end{table}
\clearpage

%%%%%%%%%%%%%%%%%%%%%%%%%%%%%%%%%%%%%%%%%%%%%%%%%%%%%%%%%%%%%%%%%%%%%%%%%%%%%%%
\subsection{Launch Site Requirements}
This is a very preliminary assessment and might change in later versions of the SED.
\begin{table}[htbp]
\centering
\caption{Launch Site Requirements}
\label{tab:launch_site_req}
\begin{tabular}{lll}
\toprule
\textbf{Category} & \textbf{Item} & \textbf{Amount} \\
\midrule
Workstation&Table & 3 \\
 &Chair & 10 \\
 &Power Supply 0-30V & 3\\
 &Power Extension Cord & 2 \\

\bottomrule
\end{tabular}
\end{table}
\clearpage

%%%%%%%%%%%%%%%%%%%%%%%%%%%%%%%%%%%%%%%%%%%%%%%%%%%%%%%%%%%%%%%%%%%%%%%%%%%%%%%
\subsection{Flight Requirements}
\begin{table}[htbp]
\centering
\caption{Flight Requirements applicable to BEXUS}
\label{tab:flight_req_bexus}
\begin{tabular}{p{5cm} p{10cm}}
\toprule
\textbf{Optimal altitude} & 22km \\
\textbf{Preferred path} & ascent \\
\textbf{Minimum float time} & not important\\
\textbf{Ground track length} & not important\\
\textbf{Light/dark condition} & no preference but should be uniform during ascent if possible\\
\bottomrule
\end{tabular}
\end{table}
%\clearpage

%%%%%%%%%%%%%%%%%%%%%%%%%%%%%%%%%%%%%%%%%%%%%%%%%%%%%%%%%%%%%%%%%%%%%%%%%%%%%%%
\subsection{Accommodation Requirements}

The experiment needs access to fresh air and the gas in- and outlets should not be covered by the canvas. They should be directed away from the vehicle and other experiments to minimize interference and contamination of the intaken air. Therefore the experiment should be accommodated in a corner of the gondola.
\clearpage

%%%%%%%%%%%%%%%%%%%%%%%%%%%%%%%%%%%%%%%%%%%%%%%%%%%%%%%%%%%%%%%%%%%%%%%%%%%%%%%
\section{Preparation and Test Activity at Esrange}
\label{sec:ESRANGE_activity}

\begin{table}[htbp]
\centering
\caption{List of planned activities}
\label{tab:preparation}
\begin{tabular}{p{1.5cm} p{2.5cm} p{4cm} p{2cm} p{1.5cm} p{2cm}}
\toprule
Time/Day & Main Task & Description & Responsible & Duration [h:m] & Comments\\
\midrule
1 & Assembly of Experiment & Assemble experiment & Antoni, Max & 03:00&\\
1 & Electronics Test & Make sure that everything is working electrically & Max, Oliver & 02:00&\\
1 & Pressurisation Test & Test pressurisation system & Leo, Antoni & 02:00 &\\
2 & Thermal Test & Test thermal system & Jonathan, Oliver & 02:00 &\\
2 & Software Test & Test Elink connection and transition from manual to autonomous mode & Abdullah & 02:00 & Make sure that the SD card is empty afterwards\\
2 & Functionality Test & Test that all sensors function properly & Jonathan, Anna & 02:00 & \\ 
\bottomrule
\end{tabular}
\end{table}

The activities in \tref{tab:preparation} are planned to take place within the first two days after arrival before the Gondola Interference Test and the Flight Compatibility Test to make sure that there is enough time to notice issues and possibly fix them. The list is preliminary at this stage and the durations are expected to change. There will be a team morning meeting to make sure that everyone is updated and delays can be taken care of quickly. It is expected that there will be some spare time on day 3 before the Gondola Interference Test.\\
The preparations involve mainly assembly and testing to verify that assembly worked well and that no component got dammaged during transport. MIRAGE doesn't have parts that need to be removed before flight or that require late access. The experiment is capable of heating sensitive components to prevent undercooling and in the case of ambient temperatures below -20°C this needs to be considered.
%\begin{table}[htbp]
%\centering
%\caption{List of planned activities}
%\label{tab:planned_activities}
%\begin{tabular}{lllll}
%\toprule
%\textbf{Time/Day} & \textbf{Main Task} & \textbf{Description} & \textbf{Responsible} & \textbf{Duration [h:m]} \\
%\midrule
%& & & & \\
%\bottomrule
%\end{tabular}
%\end{table}
\clearpage


%%%%%%%%%%%%%%%%%%%%%%%%%%%%%%%%%%%%%%%%%%%%%%%%%%%%%%%%%%%%%%%%%%%%%%%%%%%%%%%
\section{Timeline for Countdown and Flight}
\label{sec:timeline}


%\begin{table}[htbp]
%\centering
%\caption{Example of timeline of the experiment events}
%\label{tab:timeline}
%\begin{tabular}{lll}
%\toprule
%\textbf{Time [s]} & \textbf{EVENT} & \textbf{Function} \\
%\midrule
%T - .. s & POWER ON & \\
%& & \\
%& & \\
%& & \\
%\bottomrule
%\end{tabular}
%\end{table}
\clearpage



\section{Post-Flight Activities}

This will be filled out in later versions of the SED.

%Celebrate!

%%%%%%%%%%%%%%%%%%%%%%%%%%%%%%%%%%%%%%%%%%%%%%%%%%%%%%%%%%%%%%%%%%%%%%%%%%%%%%%